\subsection{CONSTRUCTION OF A FAMILY OF PSEUDOMETRICS DEFINING A UNIFORMITY}

The significance of defining a uniformity by means of a family of pseudometrics lies in the fact that all uniformities can be so obtained. Namely:

THEOREM 1. Given a uniformity \(\mathcal{U}\) on a set \(X\) , there is a family of pseudometrics on \(X\) such that the uniformity defined by this family is identical with \(\mathcal{U}\) .

For each entourage V of the uniformity U, define inductively a sequence of symmetric entourages \((\mathbf{U}_{n})\) such that \(U_{1} \subset V\) and \(\dot{U}_{n+1}^{2} \subset U_{n}\) for all \(n \geqslant 1\) . The sequence \((\mathbf{U}_{n})\) is a fundamental system of entourages of a uniformity \(U_{v}\) coarser than U; moreover, it is clear that U is the least upper bound of all the uniformities \(U_{v}\) as V runs through the filter of entourages of U. Theorem 1 is therefore a consequence of the following proposition:

PROPOSITION 2. If a uniformity U on X has a countable fundamental system of entourages, then there is a pseudometric f on X such that U is identical with the uniformity defined by f.

Let \((\mathbf{V}_n)\) be a countable fundamental system of entourages of \(\mathcal{U}\) . Define inductively a sequence \((\mathbf{U}_n)\) of symmetric entourages of \(\mathcal{U}\) such that \(\mathbf{U}_1 \subset \mathbf{V}_1\) and

\[
\mathbf {\dot {U}} _ {n + 1} ^ {3} \subset \mathbf {U} _ {n} \cap \mathbf {V} _ {n} \quad \text {   for   } \quad n \geqslant 1.
\]

Clearly \((\mathbf{U}_n)\) is another fundamental system of entourages of \(\mathcal{U}\) , and we have in particular \(\dot{\mathbf{U}}_{n+1}^3 \subset \mathbf{U}_n\) for \(n \geqslant 1\) . We define a real-valued function \(g\) on \(\mathbf{X} \times \mathbf{X}\) as follows: \(g(x, y) = 0\) if \((x, y) \in \mathbf{U}_n\) for all \(n\) ; \(g(x, y) = 2^{-k}\) if \((x, y) \in \mathbf{U}_n\) for \(1 \leqslant n \leqslant k\) , but \((x, y) \notin \mathbf{U}_{k+1}\) ; \(g(x, y) = 1\) if \((x, y) \notin \mathbf{U}_1\) . The function \(g\) is symmetric and positive, and we have \(g(x, x) = 0\) for all \(x \in \mathbf{X}\) . Put

\[
f (x, y) = \inf \sum_ {i = 0} ^ {p - 1} g \left(z _ {i}, z _ {i + 1}\right),
\]

the greatest lower bound being taken over the set of all finite sequences \((z_{i})_{0\leqslant i\leqslant p}\) ( \(p\) arbitrary) such that \(z_0 = x\) and \(z_{p} = y\) . We shall show that \(f\) is a pseudometric which satisfies the inequalities

\[
\frac {1}{2} g (x, y) \leqslant f (x, y) \leqslant g (x, y).\tag{2}
\]

It follows immediately from the definition that \(f\) is symmetric and positive and satisfies the triangle inequality. Also it is clear that \(f(x, y) \leqslant g(x, y)\) , hence \(f(x, x) = 0\) for all \(x \in \mathbf{X}\) , and therefore \(f\) is a pseudometric. To prove the left-hand half of the inequalities (2), let us show by induction on \(p\) that, for every finite sequence \((z_i)_{0 \leqslant i \leqslant p}\) of \(p + 1\) points of \(\mathbf{X}\) such that \(z_0 = x\) and \(z_p = y\) , we have

\[
\sum_ {i = 0} ^ {p - 1} g (z _ {i}, \quad z _ {i + 1}) \geqslant \frac {\mathrm{I}}{2} g (x, y).\tag{3}
\]

This is clear if \(p = \mathbf{i}\) . Put \(a = \sum_{i=0}^{p-1} g(z_i, z_{i+1})\) ; the inequality (3) is true if \(a \geqslant \mathbf{i}/2\) , because \(g(x, y) \leqslant \mathbf{i}\) . Suppose then that \(a < \mathbf{i}/2\) , and let \(h\) be the largest of the indices \(q\) such that

\[
\sum_ {i <   q} g (z _ {i}, z _ {i + 1}) \leqslant \frac {a}{2};
\]

we have then \(\sum_{i < h} g(z_i, z_{i+1}) \leqslant a/2\) and \(\sum_{i < h+1} g(z_i, z_{i+1}) > a/2\) , whence

\[
\sum_ {i > h} g (z _ {i}, z _ {i + 1}) \leqslant \frac {a}{2}.
\]

By the inductive hypothesis we have \(g(x, z_h) \leqslant a\) and \(g(z_{h+1}, y) \leqslant a\) ; on the other hand it is clear that \(g(z_h, z_{h+1}) \leqslant a\) . Let \(k\) be the smallest integer \(> 0\) such that \(2^{-k} \leqslant a\) ; then \(k \geqslant 2\) , and \((x, z_h) \in \mathbf{U}_k\) , \((z_h, z_{h+1}) \in \mathbf{U}_k\) , \((z_{h+1}, y) \in \mathbf{U}_k\) by the definition of \(g\) ; hence \((x, y) \in \mathbf{U}_k \subset \mathbf{U}_{k-1}\) , which implies that \(g(x, y) \leqslant 2^{1-k} \leqslant 2a\) .

Hence the inequalities (2) are proved; they show that, for each \(a > 0\) , the set \(\overline{f}^{1}([o, a])\) contains \(\mathbf{U}_{k}\) for each index \(k\) such that \(2^{-k} < a\) , and conversely that each \(\mathbf{U}_{k}\) contains the set \(\overline{f}^{1}([o, 2^{-k-1}])\) ; hence the sets \(\overline{f}^{1}([o, a])\) form a fundamental system of entourages of the structure \(\mathfrak{U}\) .

Remark. A uniformity \(\mathfrak{U}\) on \(\mathbf{X}\) is defined by the family \(\Phi\) of all pseudometrics on \(\mathbf{X}\) which are uniformly continuous on \(\mathbf{X} \times \mathbf{X}\) . For clearly the uniformity defined by the family \(\Phi\) is coarser than \(\mathfrak{U}\) ; conversely, Theorem 1 shows that there is a subfamily of \(\Phi\) which defines the uniformity \(\mathfrak{U}\) and therefore the uniformity defined by \(\Phi\) is finer than \(\mathfrak{U}\) .

